\section{Quotient Structure}

Starting from the $120$ flags of the dodecahedron, we form $60$ chamber states
by identifying opposite configurations along edges. Local transport operations
induce a connected tetravalent graph
\[
\Gsixty,
\]
with
\[
|V| = 60, \quad |E| = 120, \quad \deg = 4.
\]

\subsection{Two-fold quotients and covering structure}

The transport construction admits natural involutive identifications,
yielding successive $2$-fold quotients
\[
\Gsixty \longrightarrow \Gthirty \longrightarrow \Gfifteen.
\]

\begin{proposition}
Each map
\[
\Gsixty \to \Gthirty, \qquad \Gthirty \to \Gfifteen
\]
is a regular $2$-fold covering of graphs. In particular,
$\Gthirty \to \Gfifteen$ is a $\mathbb{Z}_2$-cover (equivalently, a signed
$2$-lift).
\end{proposition}

\begin{proof}
Each quotient is defined by a fixed-point-free involution, so every fiber
has size $2$.

It remains to verify that adjacency is locally preserved. Each vertex of
$\Gsixty$ (and hence of $\Gthirty$) has exactly four neighbors, corresponding
to the four transport moves on chamber states. The involutions exchange
the two representatives in a fiber without identifying distinct transport
directions. Thus the four neighbors of a vertex project to four distinct
neighbors of its image.

Since both graphs are $4$-regular, the induced map on neighborhoods is
bijective, and hence each quotient map is a covering. Regularity follows
because the involution acts transitively on each fiber, giving deck group
$\mathbb{Z}_2$.
\end{proof}

In particular, the cover $\Gthirty \to \Gfifteen$ fits into the standard
framework of voltage graphs and signed graph lifts (see
\cite{gross_tucker,zaslavsky}), and will be described explicitly by a
$\mathbb{Z}_2$-valued cocycle in Section~7.

\subsection{Identification of $\Gsixty$}

\begin{proposition}[Identification with the arc-transitive census]
The graph $\Gsixty$ constructed via dodecahedral flag transport is
isomorphic to the arc-transitive graph $\mathrm{AT4val}[60,6]$,
with identifier Graph52002 in the House of Graphs database
(\cite{potocnik_graphsym,house_of_graphs,potocnik_private}).
\end{proposition}

\begin{proof}
The graph $\Gsixty$ is connected, $4$-regular, and has $60$ vertices and
$120$ edges. Computational verification shows that its automorphism group
acts transitively on arcs and has order $480$, and that the diameter is $6$
with distance distribution $(1,4,8,16,24,6,1)$.

These invariants uniquely identify $\Gsixty$ within the census of
tetravalent arc-transitive graphs, where it is recorded as
$\mathrm{AT4val}[60,6]$. Cross-referencing with the House of Graphs
database confirms the canonical identifier Graph52002, and the
identification is confirmed in \cite{potocnik_private}.
\end{proof}
